\documentclass[10pt,a4paper]{amsart}
\usepackage[margin=19mm]{geometry}
\usepackage{amsmath,amssymb,mathtools}
\usepackage[scr=rsfs,cal=euler]{mathalfa}
\usepackage{xfrac}
\usepackage[unicode,hidelinks,bookmarks=false]{hyperref}
\newcommand{\ie}{i.e.\ }
\newcommand{\cf}{cf.\ }
\newcommand{\resp}{respectively\ }
\newcommand{\Subsection}{\subsection}
\providecommand{\nobreakdash}{\nobreak\hbox{-}}
\setlength{\emergencystretch}{2em}
\multlinegap0pt
\usepackage{bm}
\usepackage{bbm}
\usepackage{dsfont}
\usepackage{xspace}
\usepackage{cancel}
\usepackage{mathtools}
\usepackage{amsthm}
\makeatletter
\@ifundefined{theorem}{\newtheorem{theorem}{Theorem}}{}
\@ifundefined{lemma}{\newtheorem{lemma}[theorem]{Lemma}}{}
\makeatother
\title[Capacitary estimates for solutions to nonlocal Dirichlet pro]{Capacitary estimates for solutions to nonlocal Dirichlet problems}
\author{Minhyun Kim, Se-Chan Lee, Marvin Weidner}
\date{Source: arXiv:2608.09503v1 (CC BY 4.0)}
\begin{document}
\maketitle

\noindent\emph{Source and licence.} Capacitary estimates for solutions to nonlocal Dirichlet problems. Version arXiv:2608.09503v1 (CC BY 4.0), distributed under the Creative Commons Attribution 4.0 International licence, as recorded in the arXiv licence metadata for this exact version and shown on the abstract page https://arxiv.org/abs/2608.09503v1. Full rights notice: licenses/arxiv-2608.09503-CC-BY-4.0.txt.\par\smallskip

\noindent\textit{Editorial scope.} Counted unit: the Lemma whose source label is \texttt{lem-iteration} in the article (source0.tex lines 666--680, source label \texttt{lem-iteration}), delivered together with the article's own complete original proof of it (source0.tex lines 682--752, one \texttt{proof} environment of 71 lines). No mathematics is written, repaired or re-derived: statement and proof are the article's own text line for line. Required context retained uncounted: none. Every definition, display and label of those lines is retained unchanged. Omitted from this package and not redistributed: 1--665, 681--681, 753--1622. Nothing inside a retained excerpt is omitted or altered: every retained body is delivered exactly as the article's lines are written, so no omission receipt of any kind accompanies this package. Citations: the retained excerpts carry no citation command at all, so no bibliography entry of the article is transcribed and no reference is redistributed. Reference adaptation: none was needed and none was made; every reference inside the delivered unit resolves inside it, so no unresolved reference or citation is produced. Printed slips kept literal: no printed word, symbol, hypothesis or number of the counted statement or of its proof was altered. Layout and class: the article's own class is \texttt{amsart} with packages including a4wide, amssymb, mathtools, enumerate, hyperref, xcolor, esint, mathrsfs; the wrapper is the lane's pinned \texttt{amsart} editorial wrapper at 10pt on A4, which restates the theorem-like environments the retained text needs and adds the article's own macro definitions that the retained text needs (none), with extra wrapper packages (bm, bbm, dsfont, xspace, cancel, mathtools). The delivered numbering is the unit's own. Assets: no retained span contains a figure environment, a graphics-inclusion command, a PDF-page inclusion, a table or a TikZ picture; no image file is copied into this package, the package declares no asset dependency and no image file is redistributed with it. Licence: version arXiv:2608.09503v1 of this article is distributed under the Creative Commons Attribution 4.0 International licence, as recorded in the arXiv licence metadata for this exact version and shown on the abstract page https://arxiv.org/abs/2608.09503v1. A full rights notice is delivered at licenses/arxiv-2608.09503-CC-BY-4.0.txt. Characters: every retained excerpt is pure ASCII, so the delivered engine, XeLaTeX with its default Latin Modern text font, reports no missing character.
\begin{lemma}\label{lem-iteration}
Let $N \in \mathbb{N}$, $a_1>0$, $a_2>0$, $0<q<1$, and $0<\Theta_{\ast} < 1/a_1$. Suppose that a nonincreasing sequence $\{X_i\}_{i=0}^{N}$ of nonnegative real numbers and sequences $\{\theta_i\}_{i=1}^N$ and $\{h_i\}_{i=1}^N$ of nonnegative real numbers satisfy
\begin{equation*}
\theta_i\leq\Theta_\ast, \quad i=1,\dots,N,
\end{equation*}
and
\begin{equation}\label{eq-technical}
X_i \leq (1-a_1\theta_i)X_{i-1}+a_2\sum_{m=1}^{i-1}q^{i-m}(X_{m-1}-X_m) +h_i+Bq^i, \quad i=1, \dots, N,
\end{equation}
for some $B \geq 0$. Then there exist constants $c>0$ and $C>0$, depending only on $a_1$, $a_2$, $q$, and $\Theta_{\ast}$, such that for any $i=1, \dots, N$,
\begin{equation*}
X_i \leq (X_0+CB)\exp\left(-c\sum_{m=1}^{i}\theta_m\right)+C\sum_{j=1}^ih_j\exp\left(-c\sum_{m=j+1}^{i}\theta_m\right),
\end{equation*}
with the convention that $\sum_{m=i+1}^i \theta_m=0$.
\end{lemma}

\begin{proof}
Define
    \begin{equation*}
        H_0=0, \quad H_i \coloneqq \sum_{m=1}^iq^{i-m}(X_{m-1}-X_m)=X_{i-1}-X_i+qH_{i-1}, \quad i=1,\dots, N.
    \end{equation*}
Since $\{X_i\}_{i=0}^N$ is nonincreasing, we have $H_i \geq 0$ for every $i=0, \dots, N$. For $0<\delta<1$ to be chosen below, we set
    \begin{equation*}
        Y_i \coloneqq X_i+(1-\delta)H_i, \quad i=0, 1, \dots, N.
    \end{equation*}
For every $i=1, \dots, N$, using \eqref{eq-technical}, we obtain
    \begin{equation*}
    \begin{aligned}
        Y_i&=X_i+(1-\delta)(X_{i-1}-X_i)+q(1-\delta)H_{i-1}\\
        &=(1-\delta)X_{i-1}+\delta X_i+q(1-\delta)H_{i-1}\\
        &\leq (1-\delta a_1 \theta_i)X_{i-1}+\delta h_i+\delta B q^i+q(1+\delta a_2-\delta)H_{i-1}.
    \end{aligned}
    \end{equation*}


    We now choose $\delta \in (0, 1)$ sufficiently small so that
    \begin{equation*}
        q(1+\delta a_2-\delta) \leq (1-\delta a_1\Theta_{\ast})(1-\delta),
    \end{equation*}
    which is possible since $q<1$. Note that $\delta$ depends only on $a_1$, $a_2$, $q$, and $\Theta_\ast$. Since $\theta_i \leq \Theta_\ast$, this choice gives
    \begin{equation*}
        Y_i \leq (1-\delta a_1 \theta_i)Y_{i-1}+\delta h_i+\delta Bq^i.
    \end{equation*}
    Iterating this recurrence yields, for every $i=1, \dots, N$,
    \begin{equation*}
        X_i \leq Y_i \leq X_0 \prod_{m=1}^i (1-\delta a_1 \theta_m)+\delta \sum_{j=1}^i h_j \prod_{m=j+1}^i(1-\delta a_1 \theta_m)+\delta B \sum_{j=1}^i q^j \prod_{m=j+1}^i(1-\delta a_1 \theta_m),
    \end{equation*}
    with the convention that $\prod_{m=i+1}^i (1-\delta a_1 \theta_m)=1$.
    For simplicity, we set
    \begin{equation*}
        S_i \coloneqq \sum_{m=1}^i\theta_m.
    \end{equation*}
Since $1-t \leq e^{-t}$, we have
    \begin{equation*}
        \prod_{m=1}^i(1-\delta a_1 \theta_m) \leq \exp\left(-\delta a_1 S_i\right)
    \end{equation*}
    and
    \begin{equation*}
        \prod_{m=j+1}^i(1-\delta a_1 \theta_m) \leq \exp\left(-\delta a_1 (S_i-S_j)\right).
    \end{equation*}
Therefore,
    \begin{equation*}
    X_i \leq X_0 \exp\left( -\delta a_1 S_i \right) + \delta \sum_{j=1}^{i} h_j\exp\left(-\delta a_1 (S_i-S_j)\right) + \delta B \sum_{j=1}^i q^j \exp\left(-\delta a_1 (S_i-S_j)\right).
    \end{equation*}
    
    It only remains to control the term containing $B$. We claim that there exist constants $c>0$ and $C>0$, depending only on $a_1$, $a_2$, $q$, and $\Theta_\ast$, such that
    \begin{equation*}
        \sum_{j=1}^i q^j \exp\left(-\delta a_1 (S_i-S_j)\right) \leq C \exp \left(-c S_i \right).
    \end{equation*}
    Indeed, choose $j_0 \in \mathbb N \cup \{0\}$ so that 
    \begin{equation*}
       \frac{S_i}{2\Theta_{\ast}}-1 < j_0 \leq \frac{S_i}{2\Theta_{\ast}}.
    \end{equation*}
    If $j \leq j_0$, then 
    \begin{equation*}
       S_j \leq j \Theta_{\ast} \leq  j_0 \Theta_{\ast} \leq \frac{1}{2}S_i,
    \end{equation*}
    and hence
    \begin{equation*}
        \sum_{j=1}^{j_0} q^j \exp\left(-\delta a_1 (S_i-S_j)\right) \leq \exp\left(-\frac{\delta a_1}{2} S_i\right) \sum_{j=1}^{\infty}q^j=\frac{q}{1-q}\exp\left(-\frac{\delta a_1}{2}S_i\right).
    \end{equation*}
    On the other hand, we have
    \begin{equation*}
        \sum_{j=j_0+1}^i q^j \exp\left(-\delta a_1 (S_i-S_j)\right) \leq \sum_{j=j_0+1}^{\infty} q^j = \frac{q^{j_0+1}}{1-q} \leq \frac{1}{1-q} \exp\left(-\frac{|\log q|}{2\Theta_{\ast}}S_i \right).
    \end{equation*}
    A combination of these two estimates proves the claim, completing the proof.
\end{proof}

\end{document}
